\documentclass[11pt]{article}
\usepackage[utf8]{inputenc}
\usepackage{amsmath,amssymb,amsthm}
\usepackage[margin=1in]{geometry}
\usepackage{hyperref}
\usepackage{xcolor}
\usepackage{enumitem}
\usepackage{newunicodechar}
\newunicodechar{ℝ}{\ensuremath{\mathbb{R}}}
\newunicodechar{ℕ}{\ensuremath{\mathbb{N}}}

\title{Tide retrospective:\\\emph{Quartic and sextic as $k$-th specialisations (A1 follow-up refactor)}}
\author{Daniel Murfet}
\date{18 May 2026}

\begin{document}
\maketitle

\begin{abstract}
Follow-up refactor to A1 (generic even-monomial 1D template):
replaces 13 explicit-proof theorems across the quartic and sextic
files with thin specialisations of A1's generic theorems. Net
deletion 162 lines (quartic $386 \to 287$, sextic $403 \to 340$).
Single session, no roadblocks, single-attempt clean build (8298
jobs), zero \texttt{sorry} / \texttt{axiom} / \texttt{native\_decide}.
Marked as a refactor rather than a tide: no deliberation needed
(A1's deliberation handled the question of whether this refactor was
worth doing), and the work is mechanical (one bridge lemma + a
pattern of \texttt{convert ... using 3} per theorem).
\end{abstract}

\section{Setting}

A1 landed earlier today: the generic-$k$ template for
$L_k(x) = x^{2k}/(2k)!$, with nine theorems. GPT's A1 consult
recommended deferring the quartic / sextic refactor to a separate
follow-up tide, on the grounds that mixing abstraction with
downstream rewrites makes review noisier. This Tide is that
follow-up.

The value: makes A1's abstraction visible. Before, the quartic and
sextic files each carried ${\sim}25$-line proofs of theorems that
A1 now subsumes; after, those theorems are 3--5 line specialisations.
The integrability lemmas
(\texttt{quartic\_integrable\_pow*} and
\texttt{sextic\_integrable\_pow*}) are kept since A1 deferred them.

\section{The thing formalised}

Two bridge lemmas:

\begin{itemize}[leftmargin=1.5em]
\item \texttt{quarticPotential\_eq\_kthPotential : quarticPotential = kthPotential 2}
\item \texttt{sexticPotential\_eq\_kthPotential : sexticPotential = kthPotential 3}
\end{itemize}

Each proves by \texttt{funext; simp; norm\_num} (definitional up to
Nat reduction of $(2 \cdot 2)! = 24$ and $(2 \cdot 3)! = 720$).

Then 13 thin specialisations replacing the explicit proofs:

\begin{itemize}[leftmargin=1.5em]
\item Quartic ($k = 2$, 7 theorems): the half-line integral, the
  even and odd moments, the partition function (+ positivity), and
  the even / odd expected values.
\item Sextic ($k = 3$, 6 theorems): same minus the half-line
  integral, which carries an arbitrary-$m$ signature that's strictly
  more general than A1's `2j`-keyed version, so it stays as-is.
\end{itemize}

Each specialisation follows the pattern
\begin{verbatim}
rw [_Potential_eq_kthPotential]   -- when partitionFunction or
                                   -- gibbsExpectation is involved
have h := kth_* (k := 2 or 3) (by norm_num) ... ht
convert h using 3
\end{verbatim}
where \texttt{convert ... using 3} handles the Nat-cast plumbing
between the specific form ($24$, $4$ in the exponents) and the
generic form ($(2 \cdot 2)!$, $((2 \cdot 2 : \mathbb N) : \mathbb R)$).

\section{Proof strategy}

Mechanical. The bridge lemma carries the work: once
$\texttt{quarticPotential} = \texttt{kthPotential}\ 2$ is in scope,
the partitionFunction and gibbsExpectation versions specialise
directly. The moment integrals don't need the bridge because they
already use raw $\exp(-tx^4/24)$ form, which is definitionally equal
to $\exp(-tx^{2 \cdot 2}/((2 \cdot 2)!:\mathbb R))$ after Nat
reduction --- \texttt{convert} handles this transparently.

\section{Roadblocks and resolutions}

None. Each of the 13 theorems went green on first attempt with the
same \texttt{convert ... using 3} pattern.

A small mid-test diagnostic: the bridge proof initially used
\texttt{funext; simp} alone, which left
$x^4 / 24 = x^4 / \uparrow\!(\textrm{Nat.factorial}\ 4)$
unsolved. Adding \texttt{norm\_num} closes the residual Nat-cast.

\section{What was Mathlib, what was new}

\paragraph{From Mathlib.} \texttt{norm\_num}, \texttt{convert},
\texttt{funext}, \texttt{simp}.

\paragraph{From the seabed.} A1's nine theorems (consumed verbatim).

\paragraph{New here.} Two bridge lemmas (5 lines each); 13
specialisation proofs (3--5 lines each) replacing the original
${\sim}25$-line proofs.

\section{Lessons}

\paragraph{The \texttt{convert ... using 3} pattern handles
Nat-reduction noise transparently.} The specialised theorems'
statements are in the ``human'' form ($24$, $4$); the generic
theorem's RHS is in the ``Nat-cast'' form ($(2 \cdot 2)!$,
$((2 \cdot 2 : \mathbb N) : \mathbb R)$). The two are equal by Nat
computation but not syntactically. \texttt{convert ... using 3} walks
the goal three levels deep and emits the cast-equality side goals,
all of which close by definitional Nat reduction. Worth recording as
the standard specialisation idiom.

\paragraph{Bridge lemmas should be definitional whenever possible.}
The bridge identities are true by definitional unfold plus Nat
reduction; each bridge proof is one line. The alternative ---
redefining the specific potentials directly as the generic one ---
would break unfolds elsewhere in downstream files, so the
bridge-as-lemma approach is strictly less disruptive.

\section{Follow-ups}

\begin{itemize}[leftmargin=1.5em]
\item \emph{Integrability tide.} Still pending. A1 deferred
  \texttt{kth\_integrable\_pow*} and the quartic / sextic versions
  remain self-contained. Once the integrability tide lands, the
  quartic / sextic files can shrink another ${\sim}100$ lines.
\item \emph{Refactor downstream specific lemmas.} The $n = 1$
  expected-value, $n = 0$ linear case, and the affine covariance are
  kept verbatim --- they already consume the even/odd theorems
  internally, so no further gain.\footnote{Specifically
  \texttt{quartic\_expected\_value\_sq},
  \texttt{quartic\_expected\_value\_lin}, and
  \texttt{quartic\_cov\_affine}.} Skip.
\item \emph{Sextic half-line integral.} The sextic's half-line
  integral takes an arbitrary $m : \mathbb N$ (not just $2 \cdot n$).
  A1's analogue is
  keyed on $2 \cdot j$ and is therefore weaker. The sextic version
  was kept verbatim. A genuinely-arbitrary-$m$ generic version in A1
  would supersede it; deferred until needed.
\end{itemize}

\section*{Acknowledgements}

GPT-5.5~Pro's A1 consult set the structural decision to do the
refactor as a separate tide, which is what made this clean. No
fresh GPT consult was needed.

\end{document}
